\section{Introduction}\label{sec:introduction}

All Banach spaces in this paper are real. A map \(F:X\to Y\) is a
\emph{bi-Lipschitz embedding} if there are constants
\(0<c\leq C<\infty\) such that
\[
 c\|x-x'\|_X\leq\|F(x)-F(x')\|_Y
       \leq C\|x-x'\|_X\qquad(x,x'\in X).
\]
It is a \emph{bi-Lipschitz equivalence} when it is also onto \(Y\).
The distinction between these two notions is central here: the result
constructs an onto map between whole Banach spaces, while its linear
obstruction concerns a subspace.

The nonlinear geometry of Banach spaces asks how much linear structure
is retained by maps that control distances but need not respect
addition. Ribe proved that uniformly homeomorphic normed spaces have
the same finite-dimensional linear structure up to one common
distortion bound \cite[Theorem~1]{Ribe76}. Aharoni and Lindenstrauss
constructed nonseparable Banach spaces that are bi-Lipschitz equivalent
but not linearly isomorphic, and asked for separable examples
\cite[pp.~281--282]{AL78}. Ribe later produced separable uniformly
homeomorphic nonisomorphic spaces, and Aharoni and Lindenstrauss
obtained uniformly convex examples that are uniformly homeomorphic
but not Lipschitz equivalent \cite{Ribe84,AL85}.
Johnson, Lindenstrauss and Schechtman
developed both rigidity results and separable examples with different
linear structures under the weaker relation of uniform homeomorphism
\cite[Theorem~5.8 and Section~8(3)]{JLS96}. These results explain why a separable bi-Lipschitz
counterexample requires more than a uniform deformation.

The space \(c_0=c_0(\N)\) consists of real scalar sequences tending to
zero, with the supremum norm. Godefroy, Kalton and Lancien proved that a
Banach space bi-Lipschitz equivalent to a linear subspace of \(c_0\) is
linearly isomorphic to a subspace of \(c_0\). They also proved that a
Banach space bi-Lipschitz equivalent to \(c_0\) is linearly isomorphic
to \(c_0\) \cite[Theorems~2.1--2.2]{GKL00}. These theorems concern the
whole domain of an equivalence. They do not say that a Banach space
containing a bi-Lipschitz image of \(c_0\) must contain a linear copy
of \(c_0\). Kalton recorded the latter question as Problem~2, separately
from the separable Lipschitz-isomorphism question in Problem~3
\cite[pp.~21--22]{Kalton08}. H\'ajek, Johanis and Schlumprecht still described the
embedding question as apparently open in their January 2024 preprint
(arXiv:2401.00831v1), later published in 2025 \cite{HJS25}.

For Banach spaces \(X,Y\), write \(X\oplus_\infty Y\) for their product
with norm \(\|(x,y)\|=\max\{\|x\|,\|y\|\}\). The main result is the
following whole-space form of a negative answer to the \(c_0\)
embedding question.

\begin{theorem}\label{thm:main}
There is a separable real Banach space \(Z\) containing no closed linear
subspace isomorphic to \(c_0\), and an onto bi-Lipschitz map
\[
                 F:Z\oplus_\infty c_0\longrightarrow Z.
\]
\end{theorem}

The map is automatically a bijection by its lower Lipschitz bound.
The two spaces in Theorem~\ref{thm:main} cannot be linearly isomorphic:
the direct sum contains a linear copy of \(c_0\), whereas \(Z\) does
not. Restricting \(F\) to that summand gives a bi-Lipschitz embedding
of \(c_0\) into this same \(Z\). Aharoni proved that every separable
metric space embeds bi-Lipschitzly into \(c_0\) \cite{Aharoni74};
composition therefore makes this one space \(Z\) a bi-Lipschitz
universal target for separable metric spaces.

There are related separable constructions with weaker control of
small distances. Kalton constructed a separable Schur space, in which
weakly convergent sequences converge in norm, uniformly homeomorphic
to a space containing complemented \(c_0\)
\cite[Theorem~4.6 and Propositions~5.1--5.2]{Kalton04}.
The Schur property excludes a linear copy of \(c_0\).
Interleaving the two \(c_0\) summands also gives uniform \(c_0\) absorption
by that same Schur space, with a nonlinear modulus of continuity.
Sar{\i} constructed a coarse-Lipschitz embedding of \(c_0\) into a
separable dual Banach space, which contains no linear copy of \(c_0\)
\cite[Introduction and Theorem~5]{Sari26}. Coarse-Lipschitz distance
bounds allow additive errors. Theorem~\ref{thm:main} obtains global
two-sided Lipschitz bounds and an onto map.

An earlier real-separable counterexample to the Lipschitz-isomorphism
question uses a different linear obstruction: one space contains
\(c_0(\ell_2)\), the space of \(\ell_2\)-valued null sequences, and
its bi-Lipschitz equivalent partner contains no linear copy of that
space \cite[Theorem~1.1]{OpenAI2026Baseline}. Its Hilbert-block and
localized Lipschitz-free argument is independent of the construction
here. The present construction excludes ordinary \(c_0\) and yields
absorption by the same space.

\subsection*{Proof overview}

The first step is algebraic. Suppose \(E,U\) are Banach spaces and
\(P:E\oplus_\infty U\to U\) is the projection. We seek an onto
bi-Lipschitz map \(g:E\oplus_\infty U\to U\), a bounded linear map
\(Q:Z\to U\), and a Lipschitz map \(K:E\oplus_\infty U\to Z\) such that
\begin{equation}
                         QK=g-P.                         \label{eq:target-lift}
\end{equation}
Section~\ref{sec:absorption} proves that these data make
\(F(b,t)=b+K(t,Qb)\) an onto bi-Lipschitz map
\(Z\oplus_\infty E\to Z\). Thus the construction has two concrete
tasks: realize \eqref{eq:target-lift} with \(E=c_0\), and arrange that \(Z\) contains no
linear \(c_0\). The identity lifts the correction \(g-P\);
it imposes no surjectivity or right-inverse requirement on \(Q\).
This distinction matters because a Lipschitz right inverse of a linear
quotient onto a separable Banach space always yields a bounded linear
right inverse \cite[Corollary~3.2]{GK03}.

We first construct \(g\). Let \(E\) and \(U\) be countable-coordinate
copies of \(c_0\). Section~\ref{sec:weights} constructs the finite-coordinate
weights and their localized correction functions. Split the coordinates
of \(U\) into finite blocks arranged
in rows. In each block choose a vector \(v\) and a functional \(w\)
with \(wv=1\). A scalar entering from the preceding block replaces
the scalar \(wx\) of the current input block, while the rest of that
block is unchanged. Denote this frozen shift at state \(s\) by
\(S_s\). With all \(v,w\) fixed, the output readouts recover
the input by an explicit linear inverse. The weights actually depend
on the normalized input, so this frozen inverse does not by itself
invert the nonlinear map. Section~\ref{sec:shift} controls the change
of the weights to obtain global bi-Lipschitz bounds, then uses Brouwer's
theorem on finite rows. An extra scalar at the end of each row both
balances the dimensions and forces the possible output spillover to
vanish. Density and a closed-image argument then give onto-ness.

We next construct \(Z\) and lift the correction. Record each coordinate
of \((S_s-P)s\) as a scalar function of \(s\) in the unit ball of
\(E\oplus_\infty U\). These functions depend on finitely many
coordinates and have localized supports. Section~\ref{sec:tests}
enlarges them by two local operations, one making a function constant
on a small cylinder and one retaining only the change. A strict scale
budget keeps their supports localized, and a nonincreasing order of
scales controls what happens after a small local change. Tail factors
give a second way to separate supports.

Section~\ref{sec:molecules} pairs these tests with mass-zero finite
combinations of points. This uses the molecule viewpoint of Arens and
Eells \cite[Section~2]{AE56} and the Lipschitz-free framework of
Godefroy and Kalton \cite[Sections~1--2]{GK03}. Here the norm is built
from lists of selected tests whose signed sums have a common
Lipschitz bound. It combines their evaluation norms, with exponents
approaching one, in an outer square sum. The square sum prevents a
hypothetical \(c_0\) basis from escaping through different test
classes; the varying exponents then select one scalar test detecting
each basis vector. We prove directly the completion and duality
properties needed for this norm. The coordinate tests define \(Q\),
and radial extension of the point molecules defines \(K\) with
\eqref{eq:target-lift}.

Finally, a hypothetical linear \(c_0\) basis in \(Z\) has uniformly
bounded signed sums. Section~\ref{sec:no-c0} extracts one detecting
test for each vector from a single class with a common operation
budget. A test's label records its original correction coordinate.
Distinct labels can be
separated either by tails or by nested coordinate cylinders. For a
repeated label, an expansion leaves either disjoint tails or a first
small local change. In the latter case, its support shrinks and its
value becomes small enough that nested tests still have uniformly
Lipschitz signed sums. The norm then forces
the signed sums of the vectors to grow, a contradiction.
Section~\ref{sec:assembly} combines this obstruction with \eqref{eq:target-lift} and
the onto shift.
